\section{Appendix A. Mathematical Derivations}

\subsection{A.1 Galerkin Discretization of Attention}

The continuous attention operator:

\[
(\mathcal{K} Q)(s) = \int_0^1 \kappa(s, t) V(t) \, dt
\tag{A.1}
\]

Choose an orthonormal basis $\{\phi_k\}_{k=1}^r$ of $L^2[0, 1]$. The Galerkin 
approximation $\hat{Q} = \sum_k q_k \phi_k$ satisfies:

\[
\langle \mathcal{K} Q, \phi_j \rangle = \sum_k q_k \langle \mathcal{K} \phi_k, \phi_j \rangle
\tag{A.2}
\]

In matrix form, define $\Phi_{n,k} = \phi_k(t_n)$ and $Q_\Phi = \Phi^\top Q$:

\[
Q_\Phi = (\Phi^\top \mathcal{K} \Phi) Q_\Phi
\tag{A.3}
\]

where $\mathcal{K}_{ij} = \kappa(t_i, t_j)$ is the discretized kernel.

\subsection{A.2 DCT Orthonormality}

\textbf{Theorem.} The DCT-II basis vectors:

\[
\phi_k(n) = \sqrt{\frac{2 - \delta_{k0}}{N}} \cos\!\left(\frac{\pi k (2n+1)}{2N}\right)
\tag{A.4}
\]

are orthonormal: $\sum_n \phi_i(n) \phi_j(n) = \delta_{ij}$.

\textbf{Proof.} Standard DCT result (Ahmed et al., 1974). $\square$

\subsection{A.3 Approximation Error Bound}

\textbf{Theorem.} For the DCT basis, the Galerkin approximation error is:

\[
\|\mathcal{K} Q - \hat{\mathcal{K}} \hat{Q}\|_2 \leq \sqrt{\sum_{k > r} \lambda_k}
\tag{A.5}
\]

where $\lambda_k$ are the eigenvalues of $\mathcal{K}^\top \mathcal{K}$ 
sorted descending.

\textbf{Proof.} Standard result in finite element theory (Strang \& Fix, 1973, 
Theorem 5.2). $\square$

\subsection{A.4 Energy Preservation Scaling}

\textbf{Claim.} The scaling factor $\sqrt{N/r}$ ensures:

\[
\mathbb{E}\left[\|\Phi O_\Phi \sqrt{N/r}\|_2^2\right] \approx \mathbb{E}\left[\|O_\Phi\|_2^2\right] \cdot (N/r) \cdot r = N \cdot \mathbb{E}[\|O_\Phi\|^2 / r]
\tag{A.6}
\]

Under the assumption of i.i.d. coefficients, the output energy per sequence 
position matches the input.

\subsection{A.5 Partial RoPE}

For Qwen3.5 with $d_h = 256$ and $d_\text{rot} = 64$:

\[
\tilde{Q} = [\text{RoPE}(Q_{:64}), \ Q_{64:}]
\tag{A.7}
\]

Only the first 64 coordinates rotate.

\subsection{A.6 Causality Preservation Fallback}

For decoders, the projection cannot preserve causality because projected 
keys depend on all tokens. We use exact causal softmax.

\subsection{A.7 Rank Selection}

The optimal rank $r^*$ minimizes:

\[
r^* = \arg\min_r \left[ \text{FLOPs}(r) + \lambda \cdot \text{ApproxError}(r) \right]
\tag{A.8}
\]

For DCT, $\text{ApproxError}(r) \sim \exp(-c r)$ for smooth kernels, so 
$r = 64$ suffices for $N = 128$.